\section{Reproducibility}\label{sec:reproducibility}

This paper aims to score 2 on the Reproducibility Disclosure Score rubric of \citet{Khan2026Survey}: \emph{(i) code public}, \emph{(ii) data openly accessible}.

\textit{Code.} The \texttt{mr-audit} package is at \href{https://github.com/ayk5511/mr-audit}{github.com/ayk5511/mr-audit}, MIT-licensed. The empirical demonstration in \Cref{sec:demonstration} is produced by five scripts in the same repository: \texttt{studies/01\_instrument\_paper2.py} (writes the 6,825-record audit log from the Khan 2026 forecast panel), \texttt{studies/02\_use\_cases.py} (executes use cases 1--5 and writes \texttt{studies/results/use\_cases\_summary.json}), \texttt{studies/03\_scaling\_benchmark.py} (the 1{,}000- and 10{,}000-record scaling benchmark in \Cref{tab:scaling}), \texttt{studies/04\_output\_drift\_analysis.py} (the per-model output-drift analysis in \Cref{tab:output_drift}), and \texttt{studies/05\_adversarial\_robustness.py} (the controlled-tampering experiments in \Cref{tab:adversarial}). The cross-domain examples in \Cref{tab:crossdomain} are produced by \texttt{examples/sklearn\_credit\_classifier.py} and \texttt{examples/statsmodels\_sarimax\_macro.py}.

\textit{Data.} The upstream data is the Khan (2026) volatility forecast panel at \href{https://github.com/ayk5511/volatility-forecasting/blob/main/results/forecasts_5d.parquet}{ayk5511/volatility-forecasting/blob/main/results/forecasts\_5d.parquet}, CC0-licensed. The panel is itself reproducible from Yahoo Finance public data via the upstream paper's data-collection scripts in under five minutes from a clean Python environment, which makes our demonstration recursively reproducible. The instrumented audit log (\texttt{paper2\_demo.db}, 15.1 MB) is committed to the \texttt{mr-audit} repository at \texttt{studies/audit\_logs/paper2\_demo.db}, along with the regulator-export bundle (\texttt{studies/audit\_logs/regulator\_bundle.zip}, 841 KB).

\textit{Audit.} The repository includes \texttt{paper/audit.py}, which re-derives every numerical claim in this paper from the JSON outputs of the demonstration scripts. The audit script must print \texttt{ALL CHECKS PASSED} (exit 0) for the paper to be considered ship-ready. We adopt this convention from earlier papers in this series \citep{KhanVol2026, Khan2026VolEval} after a fabricated-table incident in \citet{KhanVol2026}; the audit-script convention catches that class of error mechanically. In the writing of \Cref{sec:demonstration}, the audit script caught two transcription errors (replay sample size mis-count, drift-month total off by one) before any push.

\textit{Tests.} The package ships with 54 unit tests covering the audit decorator, hash chain, replay engine, input and output drift modules, regulator-export bundle, and the SQLite and JSONL stores. The CI configuration runs them on every commit on Python 3.11 / 3.12 / 3.13 / 3.14 across Linux, macOS, and Windows.

\textit{Issues.} Corrections, extensions, and challenges to the implementation choices reported here are welcomed via the repository's issue tracker (\href{https://github.com/ayk5511/mr-audit/issues}{github.com/ayk5511/mr-audit/issues}). Issues become part of the public record and feed the next version of this paper.

\textit{Self-citation network.} This paper extends the methodological agenda of \citet{Khan2026Survey} (the Reproducibility Disclosure Score and the broader case for audit-grade artefacts) and uses the empirical artefact of \citet{KhanVol2026} (the seven-model volatility panel) as its instrumentation target. It is the technical companion to \citet{Khan2026VolEval} (the \texttt{vol-eval} package and 30-paper re-analysis): where \texttt{vol-eval} addresses the question ``did your forecasting comparison use the right significance test?'', \texttt{mr-audit} addresses the question ``can you prove tomorrow that today's prediction was made with this code, this data, and these parameters?''.
